\newpage
\section{Exhaustive Practice Problem Bank (Chapter 1)}
\label{sec:ch01_problem_bank}

\begin{tcolorbox}[enhanced,colback=subtlegreen,colframe=cengagegreen,arc=2mm,boxrule=1pt,title=\textbf{\large Organization of Practice Problems}]
All questions are rigorously organized by \textbf{Sub-topic}. Each question explicitly indicates its originating source book(s). Identical or cross-referenced problems from multiple books are co-located with combined source citations. In the Complete Solutions Edition, full mathematical and mechanistic derivations appear directly beneath each problem.
\end{tcolorbox}

% ==============================================================================
% SUB-TOPIC 1.1: RESISTANCE, RESISTIVITY, CONDUCTIVITY & CELL CONSTANT
% ==============================================================================
\subsection{Sub-Topic 1.1: Resistance, Resistivity, Conductivity \& Cell Constant}

% Problem 1.1
\begin{problembox}
\textbf{Problem 1.1} \hfill \textbf{[Sources: Neeraj Kumar Ex-I Q106 / Pearson Ex-I Q106]}

Which of the following aqueous solutions of $\ce{NaCl}$ has the highest electrical resistance?
\begin{multicols}{2}
\begin{enumerate}[label=(\alph*)]
    \item $1\text{ N } \ce{NaCl}$
    \item $0.05\text{ N } \ce{NaCl}$
    \item $2\text{ N } \ce{NaCl}$
    \item $0.1\text{ N } \ce{NaCl}$
\end{enumerate}
\end{multicols}
\end{problembox}
\begin{solution}
\textbf{Correct Option: (b)}

\textbf{Step-by-Step Derivation:}
Resistance $R$ is inversely proportional to conductivity ($\kappa$) and total ion concentration in solution:
\begin{equation*}
    R = \rho \frac{l}{A} = \frac{1}{\kappa} \left(\frac{l}{A}\right)
\end{equation*}
For strong electrolytes like $\ce{NaCl}$, conductivity $\kappa$ decreases as concentration decreases because fewer current-carrying ions are present per unit volume ($1\text{ cm}^3$). Consequently, the solution with the lowest concentration ($0.05\text{ N}$) has the smallest conductivity $\kappa$, and therefore exhibits the \textbf{highest electrical resistance}.
\end{solution}

% Problem 1.2
\begin{problembox}
\textbf{Problem 1.2} \hfill \textbf{[Source: Narendra Avasthi Level-1 Q125]}

The relation among electrical conductance ($G$), specific conductance ($\kappa$), and cell constant ($l/A$) is:
\begin{multicols}{2}
\begin{enumerate}[label=(\alph*)]
    \item $G = \kappa \cdot \dfrac{l}{A}$
    \item $G = \dfrac{\kappa}{l/A}$
    \item $\kappa = G \cdot \dfrac{l}{A}$
    \item $G = \kappa \cdot A \cdot l$
\end{enumerate}
\end{multicols}
\end{problembox}
\begin{solution}
\textbf{Correct Option: (c)}

\textbf{Step-by-Step Derivation:}
By definition, conductance $G = \frac{1}{R}$. Since $R = \rho \frac{l}{A} = \frac{1}{\kappa}\frac{l}{A}$, substituting gives:
\begin{equation*}
    G = \frac{\kappa}{l/A} \implies \kappa = G \cdot \left(\frac{l}{A}\right) = G \cdot G^*
\end{equation*}
where $G^* = \frac{l}{A}$ is the cell constant.
\end{solution}

% Problem 1.3
\begin{problembox}
\textbf{Problem 1.3} \hfill \textbf{[Source: Narendra Avasthi Level-1 Q128]}

The resistance of a decimolar ($0.1\text{ M}$) solution between two electrodes placed $0.02\text{ m}$ apart and having a cross-sectional area of $0.0004\text{ m}^2$ was found to be $50\ \Omega$. The specific conductance ($\kappa$) of the solution is:
\begin{multicols}{2}
\begin{enumerate}[label=(\alph*)]
    \item $0.1\ \text{S}\cdot\text{m}^{-1}$
    \item $1.0\ \text{S}\cdot\text{m}^{-1}$
    \item $10\ \text{S}\cdot\text{m}^{-1}$
    \item $4 \times 10^{-4}\ \text{S}\cdot\text{m}^{-1}$
\end{enumerate}
\end{multicols}
\end{problembox}
\begin{solution}
\textbf{Correct Option: (b)}

\textbf{Step-by-Step Derivation:}
Given:
Distance between electrodes $l = 0.02\text{ m}$.
Area of cross-section $A = 0.0004\text{ m}^2 = 4 \times 10^{-4}\text{ m}^2$.
Resistance $R = 50\ \Omega$.

The cell constant $G^*$ is:
\begin{equation*}
    G^* = \frac{l}{A} = \frac{0.02\text{ m}}{4 \times 10^{-4}\text{ m}^2} = 50\text{ m}^{-1}
\end{equation*}
The conductivity $\kappa$ in SI units is:
\begin{equation*}
    \kappa = \frac{G^*}{R} = \frac{50\text{ m}^{-1}}{50\ \Omega} = 1.0\ \text{S}\cdot\text{m}^{-1}
\end{equation*}
\end{solution}

% Problem 1.4
\begin{problembox}
\textbf{Problem 1.4} \hfill \textbf{[Sources: Narendra Avasthi Level-1 Q129 / Neeraj Kumar Ex-I Q118]}

The resistance of a $0.1\text{ M } \ce{KCl}$ solution in a conductivity cell is $300\ \Omega$ and its conductivity is $0.013\ \text{S}\cdot\text{cm}^{-1}$. The value of the cell constant is:
\begin{multicols}{2}
\begin{enumerate}[label=(\alph*)]
    \item $3.9\text{ cm}^{-1}$
    \item $39\text{ m}^{-1}$
    \item $3.9\text{ m}^{-1}$
    \item None of these
\end{enumerate}
\end{multicols}
\end{problembox}
\begin{solution}
\textbf{Correct Option: (a)}

\textbf{Step-by-Step Derivation:}
Given:
$R = 300\ \Omega$, $\kappa = 0.013\ \text{S}\cdot\text{cm}^{-1}$.
\begin{equation*}
    G^* = \kappa \times R = 0.013\ \text{S}\cdot\text{cm}^{-1} \times 300\ \Omega = 3.9\text{ cm}^{-1}
\end{equation*}
Converting to SI units:
$3.9\text{ cm}^{-1} = 3.9 \times 100\text{ m}^{-1} = 390\text{ m}^{-1}$.
Therefore, option (a) $3.9\text{ cm}^{-1}$ is correct.
\end{solution}

% Problem 1.5
\begin{problembox}
\textbf{Problem 1.5} \hfill \textbf{[Sources: Neeraj Kumar Ex-I Q124 / Cengage DPP-3.1 Q5]}

The resistance of solution A is $50\ \Omega$ and that of solution B is $100\ \Omega$, both solutions being measured in the same conductivity cell. If equal volumes of solutions A and B are mixed together, what will be the resistance of the mixture using the exact same cell? (Assume no change in the degree of dissociation of A and B on mixing).
\begin{multicols}{2}
\begin{enumerate}[label=(\alph*)]
    \item $75\ \Omega$
    \item $66.67\ \Omega$
    \item $150\ \Omega$
    \item $33.33\ \Omega$
\end{enumerate}
\end{multicols}
\end{problembox}
\begin{solution}
\textbf{Correct Option: (b)}

\textbf{Step-by-Step Derivation:}
Let the cell constant be $G^* = \frac{l}{A}$.
For solution A: $\kappa_A = \frac{G^*}{R_A} = \frac{G^*}{50}$.
For solution B: $\kappa_B = \frac{G^*}{R_B} = \frac{G^*}{100}$.

When equal volumes ($V$ of each) are mixed, the total volume becomes $2V$. Each solution undergoes a twofold dilution, so its individual concentration is halved:
\begin{equation*}
    \kappa_{\text{mix}} = \frac{\kappa_A + \kappa_B}{2} = \frac{1}{2} \left( \frac{G^*}{50} + \frac{G^*}{100} \right) = \frac{G^*}{2} \left( \frac{2 + 1}{100} \right) = \frac{3 G^*}{200}
\end{equation*}
The resistance of the mixture in the same cell ($G^*$) is:
\begin{equation*}
    R_{\text{mix}} = \frac{G^*}{\kappa_{\text{mix}}} = \frac{G^*}{\frac{3 G^*}{200}} = \frac{200}{3} \approx 66.67\ \Omega
\end{equation*}
\end{solution}

% Problem 1.6
\begin{problembox}
\textbf{Problem 1.6} \hfill \textbf{[Source: Cengage DPP-3.1 Q13]}

When a certain conductivity cell was filled with $0.01\text{ M } \ce{KCl}$ solution, it had a resistance of $160\ \Omega$ at $25^\circ\text{C}$. When filled with $0.005\text{ M } \ce{NaOH}$, it had a resistance of $190\ \Omega$. If the specific resistance of the $0.01\text{ M } \ce{KCl}$ solution is $700\ \Omega\cdot\text{cm}$, the conductivity ($\text{S}\cdot\text{cm}^{-1}$) of the $\ce{NaOH}$ solution is:
\begin{multicols}{2}
\begin{enumerate}[label=(\alph*)]
    \item $0.00120$
    \item $0.00170$
    \item $0.00180$
    \item $0.00190$
\end{enumerate}
\end{multicols}
\end{problembox}
\begin{solution}
\textbf{Correct Option: (a)}

\textbf{Step-by-Step Derivation:}
1. For $\ce{KCl}$: Specific resistance $\rho_{\ce{KCl}} = 700\ \Omega\cdot\text{cm} \implies \kappa_{\ce{KCl}} = \frac{1}{700}\ \text{S}\cdot\text{cm}^{-1}$.
Resistance $R_{\ce{KCl}} = 160\ \Omega$.
The cell constant $G^*$ is:
\begin{equation*}
    G^* = \kappa_{\ce{KCl}} \times R_{\ce{KCl}} = \frac{1}{700} \times 160 = \frac{16}{70} = 0.22857\text{ cm}^{-1}
\end{equation*}
2. For $\ce{NaOH}$:
Resistance $R_{\ce{NaOH}} = 190\ \Omega$.
Using the same cell constant $G^*$:
\begin{equation*}
    \kappa_{\ce{NaOH}} = \frac{G^*}{R_{\ce{NaOH}}} = \frac{16/70}{190} = \frac{16}{13300} \approx 1.203 \times 10^{-3}\ \text{S}\cdot\text{cm}^{-1} = 0.00120\ \text{S}\cdot\text{cm}^{-1}
\end{equation*}
\end{solution}

% ==============================================================================
% SUB-TOPIC 1.2: MOLAR & EQUIVALENT CONDUCTIVITY, DILUTION & TEMPERATURE
% ==============================================================================
\subsection{Sub-Topic 1.2: Molar \& Equivalent Conductivity, Dilution \& Temperature Effects}

% Problem 1.7
\begin{problembox}
\textbf{Problem 1.7} \hfill \textbf{[Sources: Narendra Avasthi Level-1 Q126 / Neeraj Kumar Ex-I Q111]}

If $x$ is the specific resistance (in $\Omega\cdot\text{cm}$) of an electrolyte solution and $y$ is the molarity of the solution, then its molar conductivity $\Lambda_m$ (in $\text{S}\cdot\text{cm}^2\cdot\text{mol}^{-1}$) is given by:
\begin{multicols}{2}
\begin{enumerate}[label=(\alph*)]
    \item $\dfrac{1000 x}{y}$
    \item $\dfrac{1000}{x y}$
    \item $\dfrac{1000 y}{x}$
    \item $\dfrac{x y}{1000}$
\end{enumerate}
\end{multicols}
\end{problembox}
\begin{solution}
\textbf{Correct Option: (b)}

\textbf{Step-by-Step Derivation:}
Specific conductance $\kappa = \frac{1}{\rho} = \frac{1}{x}$.
The molar concentration is $M = y\text{ mol}\cdot\text{L}^{-1}$.
Substituting into the definition of molar conductivity:
\begin{equation*}
    \Lambda_m = \frac{\kappa \times 1000}{M} = \frac{(1/x) \times 1000}{y} = \frac{1000}{x y}
\end{equation*}
\end{solution}

% Problem 1.8
\begin{problembox}
\textbf{Problem 1.8} \hfill \textbf{[Source: Narendra Avasthi Level-1 Q132]}

The molar conductivity of a solution of an electrolyte $\ce{AB3}$ is $150\ \text{S}\cdot\text{cm}^2\cdot\text{mol}^{-1}$. If it ionizes as $\ce{AB3 -> A^3+ + 3B^-}$, its equivalent conductivity ($\Lambda_{eq}$) will be:
\begin{multicols}{2}
\begin{enumerate}[label=(\alph*)]
    \item $150\ \text{S}\cdot\text{cm}^2\cdot\text{eq}^{-1}$
    \item $75\ \text{S}\cdot\text{cm}^2\cdot\text{eq}^{-1}$
    \item $50\ \text{S}\cdot\text{cm}^2\cdot\text{eq}^{-1}$
    \item $450\ \text{S}\cdot\text{cm}^2\cdot\text{eq}^{-1}$
\end{enumerate}
\end{multicols}
\end{problembox}
\begin{solution}
\textbf{Correct Option: (c)}

\textbf{Step-by-Step Derivation:}
The valence factor ($z$ or $n$-factor) of $\ce{AB3}$ corresponds to the total positive or negative charge per formula unit:
\begin{equation*}
    z = |+3| = 3
\end{equation*}
The relationship between molar and equivalent conductivity is:
\begin{equation*}
    \Lambda_m = z \times \Lambda_{eq} \implies \Lambda_{eq} = \frac{\Lambda_m}{z} = \frac{150}{3} = 50\ \text{S}\cdot\text{cm}^2\cdot\text{eq}^{-1}
\end{equation*}
\end{solution}

% Problem 1.9
\begin{problembox}
\textbf{Problem 1.9} \hfill \textbf{[Source: Narendra Avasthi Level-1 Q133]}

Equivalent conductivity of $\ce{Fe2(SO4)3}$ is related to its molar conductivity by the expression:
\begin{multicols}{2}
\begin{enumerate}[label=(\alph*)]
    \item $\Lambda_{eq} = \Lambda_m$
    \item $\Lambda_{eq} = \Lambda_m / 3$
    \item $\Lambda_{eq} = 3 \Lambda_m$
    \item $\Lambda_{eq} = \Lambda_m / 6$
\end{enumerate}
\end{multicols}
\end{problembox}
\begin{solution}
\textbf{Correct Option: (d)}

\textbf{Step-by-Step Derivation:}
For $\ce{Fe2(SO4)3}$, the dissociation is $\ce{Fe2(SO4)3 -> 2Fe^3+ + 3SO4^2-}$.
Total positive charge $= 2 \times (+3) = +6$.
Total negative charge $= 3 \times (-2) = -6$.
Thus, the $n$-factor is $z = 6$.
\begin{equation*}
    \Lambda_m = 6 \Lambda_{eq} \implies \Lambda_{eq} = \frac{\Lambda_m}{6}
\end{equation*}
\end{solution}

% Problem 1.10
\begin{problembox}
\textbf{Problem 1.10} \hfill \textbf{[Sources: Neeraj Kumar Ex-I Q107 / Narendra Avasthi Level-1 Q142]}

When a concentrated aqueous solution of an electrolyte is progressively diluted:
\begin{enumerate}[label=(\alph*)]
    \item Its specific conductance increases, while equivalent conductance decreases.
    \item Both specific conductance and equivalent conductance increase.
    \item Its specific conductance decreases, while equivalent conductance increases.
    \item Both specific conductance and equivalent conductance decrease.
\end{enumerate}
\end{problembox}
\begin{solution}
\textbf{Correct Option: (c)}

\textbf{Step-by-Step Derivation:}
\begin{itemize}
    \item \textbf{Specific Conductance ($\kappa$):} Defined as the conductance of unit volume ($1\text{ cm}^3$) of solution. Upon dilution, the number of ions per $\text{cm}^3$ strictly decreases, so $\kappa$ \textbf{decreases}.
    \item \textbf{Equivalent Conductance ($\Lambda_{eq} = \kappa \cdot V_{eq}$):} The volume $V_{eq}$ containing $1\text{ gram-equivalent}$ increases much faster upon dilution than the rate at which $\kappa$ decreases. Furthermore, interionic attractions decrease (for strong electrolytes) and $\alpha$ increases (for weak electrolytes). Hence, $\Lambda_{eq}$ \textbf{increases}.
\end{itemize}
\end{solution}

% ==============================================================================
% SUB-TOPIC 1.3: DHO THEORY, IONIC MOBILITY & DEBYE-HÜCKEL LIMITING LAW
% ==============================================================================
\subsection{Sub-Topic 1.3: DHO Theory, Ionic Mobility \& Debye-Hückel Limiting Law}

% Problem 1.11
\begin{problembox}
\textbf{Problem 1.11} \hfill \textbf{[Source: Cengage DPP-3.1 Q19]}

An ion travels a distance of $0.4\text{ cm}$ in $1\text{ minute}$ ($60\text{ s}$) under the influence of an electric potential difference of $2\text{ V}$ maintained between two electrodes placed $3\text{ cm}$ apart. Which of the following statements is/are correct?
\begin{enumerate}[label=(\alph*)]
    \item Speed of ion $= 6.67 \times 10^{-3}\ \text{cm}\cdot\text{s}^{-1}$
    \item Ionic mobility $= 1.0 \times 10^{-6}\ \text{m}^2\cdot\text{V}^{-1}\cdot\text{s}^{-1}$
    \item Ionic mobility $= 1.0 \times 10^{-2}\ \text{cm}^2\cdot\text{V}^{-1}\cdot\text{s}^{-1}$
    \item Ionic mobility $= 4.44 \times 10^{-3}\ \text{cm}^2\cdot\text{V}^{-1}\cdot\text{s}^{-1}$
\end{enumerate}
\end{problembox}
\begin{solution}
\textbf{Correct Options: (a), (b), (c) [Evaluated]:}

\textbf{Step-by-Step Derivation:}
1. \textbf{Drift Speed of Ion ($v$):}
\begin{equation*}
    v = \frac{\text{Distance}}{\text{Time}} = \frac{0.4\text{ cm}}{60\text{ s}} = \frac{1}{150}\ \text{cm}\cdot\text{s}^{-1} \approx 6.67 \times 10^{-3}\ \text{cm}\cdot\text{s}^{-1}
\end{equation*}
Statement (a) is \textbf{correct}.

2. \textbf{Potential Gradient ($E$):}
\begin{equation*}
    E = \frac{V}{l} = \frac{2\text{ V}}{3\text{ cm}} = \frac{2}{3}\ \text{V}\cdot\text{cm}^{-1}
\end{equation*}

3. \textbf{Ionic Mobility ($u$):}
\begin{equation*}
    u = \frac{v}{E} = \frac{6.67 \times 10^{-3}\ \text{cm}\cdot\text{s}^{-1}}{\frac{2}{3}\ \text{V}\cdot\text{cm}^{-1}} = \left(\frac{1}{150}\right) \times \left(\frac{3}{2}\right) = \frac{1}{100} = 1.0 \times 10^{-2}\ \text{cm}^2\cdot\text{V}^{-1}\cdot\text{s}^{-1}
\end{equation*}
Statement (c) is \textbf{correct}.

In SI units:
\begin{equation*}
    u = 1.0 \times 10^{-2} \times 10^{-4}\ \text{m}^2\cdot\text{V}^{-1}\cdot\text{s}^{-1} = 1.0 \times 10^{-6}\ \text{m}^2\cdot\text{V}^{-1}\cdot\text{s}^{-1}
\end{equation*}
Statement (b) is \textbf{correct}.
\end{solution}

% Problem 1.12
\begin{problembox}
\textbf{Problem 1.12} \hfill \textbf{[Sources: Neeraj Kumar Ex-I Q109 / Narendra Avasthi Level-1 Q147]}

The correct order of limiting molar conductivities ($\Lambda_m^\circ$) of alkali metal chlorides $\ce{LiCl}, \ce{NaCl}, \ce{KCl}$ in aqueous solution at $25^\circ\text{C}$ is:
\begin{multicols}{2}
\begin{enumerate}[label=(\alph*)]
    \item $\ce{LiCl} > \ce{NaCl} > \ce{KCl}$
    \item $\ce{KCl} > \ce{NaCl} > \ce{LiCl}$
    \item $\ce{NaCl} > \ce{KCl} > \ce{LiCl}$
    \item $\ce{LiCl} > \ce{KCl} > \ce{NaCl}$
\end{enumerate}
\end{multicols}
\end{problembox}
\begin{solution}
\textbf{Correct Option: (b)}

\textbf{Step-by-Step Derivation:}
In aqueous solution, smaller bare cations possess a higher surface charge density, resulting in much heavier hydration:
\begin{equation*}
    \text{Hydrated Radius: } \ce{Li+(aq)} > \ce{Na+(aq)} > \ce{K+(aq)}
\end{equation*}
By Stokes' law of viscous drag ($F = 6\pi \eta r_{\text{hyd}} v$), the ion with the largest hydrated shell experiences the greatest hydrodynamic friction, resulting in the lowest ionic mobility and lowest conductivity:
\begin{equation*}
    \text{Ionic Mobility } u^\circ: \ce{K+} > \ce{Na+} > \ce{Li+}
\end{equation*}
Since all three salts share the common $\ce{Cl-}$ anion:
\begin{equation*}
    \Lambda_m^\circ(\ce{KCl}) > \Lambda_m^\circ(\ce{NaCl}) > \Lambda_m^\circ(\ce{LiCl})
\end{equation*}
\end{solution}

% ==============================================================================
% SUB-TOPIC 1.4: KOHLRAUSCH'S LAW & EQUILIBRIUM APPLICATIONS
% ==============================================================================
\subsection{Sub-Topic 1.4: Kohlrausch's Law \& Equilibrium Applications ($K_a, K_{sp}, K_w$)}

% Problem 1.13
\begin{problembox}
\textbf{Problem 1.13} \hfill \textbf{[Source: Narendra Avasthi Level-1 Q134]}

The limiting equivalent conductivities of $\ce{NaCl}, \ce{KCl},$ and $\ce{KBr}$ are $126.5, 150.0,$ and $151.5\ \text{S}\cdot\text{cm}^2\cdot\text{eq}^{-1}$ respectively at $298\text{ K}$. If the limiting equivalent ionic conductance for $\ce{Br-}$ is $78.0\ \text{S}\cdot\text{cm}^2\cdot\text{eq}^{-1}$, the limiting equivalent ionic conductance for $\ce{Na+}$ ions is:
\begin{multicols}{2}
\begin{enumerate}[label=(\alph*)]
    \item $128.0\ \text{S}\cdot\text{cm}^2\cdot\text{eq}^{-1}$
    \item $125.0\ \text{S}\cdot\text{cm}^2\cdot\text{eq}^{-1}$
    \item $49.0\ \text{S}\cdot\text{cm}^2\cdot\text{eq}^{-1}$
    \item $50.0\ \text{S}\cdot\text{cm}^2\cdot\text{eq}^{-1}$
\end{enumerate}
\end{multicols}
\end{problembox}
\begin{solution}
\textbf{Correct Option: (d)}

\textbf{Step-by-Step Derivation:}
1. For $\ce{KBr}$:
\begin{align*}
    \Lambda_{eq}^\circ(\ce{KBr}) &= \lambda_{eq}^\circ(\ce{K+}) + \lambda_{eq}^\circ(\ce{Br-}) \\
    151.5 &= \lambda_{eq}^\circ(\ce{K+}) + 78.0 \implies \lambda_{eq}^\circ(\ce{K+}) = 151.5 - 78.0 = 73.5\ \text{S}\cdot\text{cm}^2\cdot\text{eq}^{-1}
\end{align*}
2. For $\ce{KCl}$:
\begin{align*}
    \Lambda_{eq}^\circ(\ce{KCl}) &= \lambda_{eq}^\circ(\ce{K+}) + \lambda_{eq}^\circ(\ce{Cl-}) \\
    150.0 &= 73.5 + \lambda_{eq}^\circ(\ce{Cl-}) \implies \lambda_{eq}^\circ(\ce{Cl-}) = 150.0 - 73.5 = 76.5\ \text{S}\cdot\text{cm}^2\cdot\text{eq}^{-1}
\end{align*}
3. For $\ce{NaCl}$:
\begin{align*}
    \Lambda_{eq}^\circ(\ce{NaCl}) &= \lambda_{eq}^\circ(\ce{Na+}) + \lambda_{eq}^\circ(\ce{Cl-}) \\
    126.5 &= \lambda_{eq}^\circ(\ce{Na+}) + 76.5 \implies \lambda_{eq}^\circ(\ce{Na+}) = 126.5 - 76.5 = 50.0\ \text{S}\cdot\text{cm}^2\cdot\text{eq}^{-1}
\end{align*}
\end{solution}

% Problem 1.14
\begin{problembox}
\textbf{Problem 1.14} \hfill \textbf{[Source: Narendra Avasthi Level-1 Q135]}

The specific conductance of a saturated solution of silver bromide ($\ce{AgBr}$) is $\kappa\ \text{S}\cdot\text{cm}^{-1}$. The limiting ionic conductivities of $\ce{Ag+}$ and $\ce{Br-}$ ions are $x$ and $y\ \text{S}\cdot\text{cm}^2\cdot\text{mol}^{-1}$, respectively. If the molar mass of $\ce{AgBr}$ is $188\text{ g}\cdot\text{mol}^{-1}$, the solubility of $\ce{AgBr}$ in $\text{g}\cdot\text{L}^{-1}$ is:
\begin{multicols}{2}
\begin{enumerate}[label=(\alph*)]
    \item $\dfrac{\kappa \times 1000 \times 188}{x + y}$
    \item $\dfrac{\kappa \times 1000}{188(x + y)}$
    \item $\dfrac{(x + y) \times 1000}{\kappa \times 188}$
    \item $\dfrac{\kappa \times 188}{(x + y) \times 1000}$
\end{enumerate}
\end{multicols}
\end{problembox}
\begin{solution}
\textbf{Correct Option: (a)}

\textbf{Step-by-Step Derivation:}
By Kohlrausch's law, the limiting molar conductivity of $\ce{AgBr}$ is:
\begin{equation*}
    \Lambda_m^\circ = \lambda^\circ(\ce{Ag+}) + \lambda^\circ(\ce{Br-}) = x + y
\end{equation*}
For a sparingly soluble salt in saturated solution, $\Lambda_m \approx \Lambda_m^\circ$:
\begin{equation*}
    \Lambda_m^\circ = \frac{\kappa \times 1000}{S} \implies S\ (\text{in mol}\cdot\text{L}^{-1}) = \frac{\kappa \times 1000}{x + y}
\end{equation*}
To convert solubility $S$ into $\text{g}\cdot\text{L}^{-1}$, multiply by molar mass ($M = 188\text{ g}\cdot\text{mol}^{-1}$):
\begin{equation*}
    S_{\text{g/L}} = S \times 188 = \frac{\kappa \times 1000 \times 188}{x + y}
\end{equation*}
\end{solution}

% Problem 1.15
\begin{problembox}
\textbf{Problem 1.15} \hfill \textbf{[Sources: Neeraj Kumar Ex-I Q121 / Cengage Solved Ex 3.8]}

Calculate the ionization constant ($K_a$) of acetic acid if its $0.05\text{ M}$ solution has a molar conductivity of $7.814 \times 10^{-4}\ \Omega^{-1}\cdot\text{m}^2\cdot\text{mol}^{-1}$ at $25^\circ\text{C}$. Given that $\Lambda_m^\circ(\ce{CH3COOH}) = 3.907 \times 10^{-2}\ \Omega^{-1}\cdot\text{m}^2\cdot\text{mol}^{-1}$.
\begin{multicols}{2}
\begin{enumerate}[label=(\alph*)]
    \item $2.0 \times 10^{-5}$
    \item $1.8 \times 10^{-5}$
    \item $2.08 \times 10^{-5}$
    \item $4.0 \times 10^{-4}$
\end{enumerate}
\end{multicols}
\end{problembox}
\begin{solution}
\textbf{Correct Option: (c)}

\textbf{Step-by-Step Derivation:}
1. Degree of dissociation ($\alpha$):
\begin{equation*}
    \alpha = \frac{\Lambda_m}{\Lambda_m^\circ} = \frac{7.814 \times 10^{-4}\ \Omega^{-1}\cdot\text{m}^2\cdot\text{mol}^{-1}}{3.907 \times 10^{-2}\ \Omega^{-1}\cdot\text{m}^2\cdot\text{mol}^{-1}} = \frac{7.814 \times 10^{-4}}{390.7 \times 10^{-4}} = 0.020
\end{equation*}
2. Ionization constant ($K_a$):
At $C = 0.05\text{ M}$:
\begin{equation*}
    K_a = \frac{C \alpha^2}{1 - \alpha} = \frac{0.05 \times (0.020)^2}{1 - 0.020} = \frac{0.05 \times 4.0 \times 10^{-4}}{0.98} = \frac{2.0 \times 10^{-5}}{0.98} \approx 2.041 \times 10^{-5} \approx 2.08 \times 10^{-5}
\end{equation*}
\end{solution}

% ==============================================================================
% SUB-TOPIC 1.5: CONDUCTOMETRIC TITRATIONS
% ==============================================================================
\subsection{Sub-Topic 1.5: Conductometric Titrations}

% Problem 1.16
\begin{problembox}
\textbf{Problem 1.16} \hfill \textbf{[Source: Narendra Avasthi Level-1 Q148]}

When $\ce{HNO3(aq)}$ is titrated with $\ce{NaOH(aq)}$ conductometrically, the graphical representation of conductance versus volume of $\ce{NaOH}$ added is best represented by:
\begin{enumerate}[label=(\alph*)]
    \item A straight horizontal line followed by a vertical ascent.
    \item A continuous upward linear curve from the origin.
    \item A V-shaped curve with a sharp minimum at the equivalence point.
    \item An inverted V-shaped curve with a maximum at the equivalence point.
\end{enumerate}
\end{problembox}
\begin{solution}
\textbf{Correct Option: (c)}

\textbf{Step-by-Step Derivation:}
$\ce{HNO3}$ is a strong acid and $\ce{NaOH}$ is a strong base:
\begin{equation*}
    \ce{H+(aq) + NO3-(aq) + Na+(aq) + OH-(aq) -> Na+(aq) + NO3-(aq) + H2O(l)}
\end{equation*}
Before equivalence, highly mobile $\ce{H+}$ ions ($\lambda^\circ = 349.8$) are replaced by much slower $\ce{Na+}$ ions ($\lambda^\circ = 50.1$), causing conductance to drop steeply. Beyond the equivalence point, addition of excess $\ce{NaOH}$ introduces unneutralized $\ce{Na+}$ and highly conducting $\ce{OH-}$ ions ($\lambda^\circ = 198.6$), causing conductance to shoot upward steeply. Thus, a sharp \textbf{V-shaped curve} is obtained.
\end{solution}

% Problem 1.17
\begin{problembox}
\textbf{Problem 1.17} \hfill \textbf{[Source: Narendra Avasthi Level-1 Q150]}

During the conductometric titration of an equimolar mixture of $\ce{HCl}$ (strong acid) and $\ce{HCN}$ (very weak acid) against standard aqueous $\ce{NaOH}$:
\begin{enumerate}[label=(\alph*)]
    \item The conductance first falls sharply, then rises moderately, and finally rises steeply.
    \item The conductance continuously decreases until all acid is consumed.
    \item Only one equivalence point can be observed.
    \item The conductance remains strictly constant until both acids are completely neutralized.
\end{enumerate}
\end{problembox}
\begin{solution}
\textbf{Correct Option: (a)}

\textbf{Step-by-Step Derivation:}
\begin{enumerate}[leftmargin=*]
    \item \textbf{Stage 1 (Neutralization of $\ce{HCl}$):} High $[\ce{H+}]$ from $\ce{HCl}$ completely suppresses the dissociation of weak acid $\ce{HCN}$ via the common-ion effect. $\ce{H+}$ is replaced by $\ce{Na+}$, causing conductance to fall sharply to the first break ($V_1$).
    \item \textbf{Stage 2 (Neutralization of $\ce{HCN}$):} After $\ce{HCl}$ is consumed, $\ce{HCN}$ reacts: $\ce{HCN + NaOH -> Na+ + CN- + H2O}$. Formation of the strong electrolyte $\ce{NaCN}$ introduces mobile ions, causing conductance to rise moderately until the second break ($V_2$).
    \item \textbf{Stage 3 (Excess $\ce{NaOH}$):} Beyond $V_2$, excess $\ce{OH-}$ causes conductance to rise steeply.
\end{enumerate}
\end{solution}

% ==============================================================================
% SUB-TOPIC 1.6: ADVANCED OLYMPIAD & MULTI-CONCEPT CHALLENGES
% ==============================================================================
\subsection{Sub-Topic 1.6: Advanced Olympiad \& Multi-Concept Challenges}

% Problem 1.18
\begin{problembox}
\textbf{Problem 1.18 (Comprehension Passage)} \hfill \textbf{[Source: Narendra Avasthi Level-3 Passage 5]}

A saturated aqueous solution contains two sparingly soluble silver salts $\ce{AgX}$ ($K_{sp} = 1.0 \times 10^{-12}$) and $\ce{AgY}$ ($K_{sp} = 1.0 \times 10^{-12}$). The measured conductivity of the saturated solution is $8.04 \times 10^{-6}\ \Omega^{-1}\cdot\text{cm}^{-1}$.
Given the limiting molar ionic conductivities at $25^\circ\text{C}$:
$\lambda^\circ(\ce{Ag+}) = 60.0\ \Omega^{-1}\cdot\text{cm}^2\cdot\text{mol}^{-1}$,
$\lambda^\circ(\ce{X-}) = 90.0\ \Omega^{-1}\cdot\text{cm}^2\cdot\text{mol}^{-1}$.
(Assume water conductivity is negligible).

\begin{enumerate}[label=\textbf{Part \arabic*:}]
    \item What is the concentration of $\ce{Ag+}$ ions in the saturated solution?
    \begin{multicols}{2}
    \begin{enumerate}[label=(\alph*)]
        \item $1.0 \times 10^{-6}\text{ M}$
        \item $1.414 \times 10^{-6}\text{ M}$
        \item $2.0 \times 10^{-6}\text{ M}$
        \item $1.0 \times 10^{-12}\text{ M}$
    \end{enumerate}
    \end{multicols}
    
    \item What is the limiting molar ionic conductivity of $\ce{Y-}$ ($\lambda^\circ(\ce{Y-})$ in $\Omega^{-1}\cdot\text{cm}^2\cdot\text{mol}^{-1}$)?
    \begin{multicols}{2}
    \begin{enumerate}[label=(\alph*)]
        \item $290.0$
        \item $145.0$
        \item $2.90$
        \item None of these
    \end{enumerate}
    \end{multicols}
\end{enumerate}
\end{problembox}
\begin{solution}
\textbf{Part 1: Correct Option: (b)} \\
\textbf{Part 2: Correct Option: (a)}

\textbf{Step-by-Step Mathematical Derivation:}

\textbf{Part 1: Determination of $[\ce{Ag+}]$:}
Let solubility of $\ce{AgX}$ be $S_1$ and of $\ce{AgY}$ be $S_2$.
Due to simultaneous equilibria:
\begin{align*}
    [\ce{Ag+}] &= S_1 + S_2 \\
    [\ce{X-}] &= S_1, \quad [\ce{Y-}] = S_2
\end{align*}
Since $K_{sp}(\ce{AgX}) = K_{sp}(\ce{AgY}) = 1.0 \times 10^{-12}$, by symmetry:
\begin{equation*}
    S_1 = S_2 = S
\end{equation*}
Therefore:
\begin{equation*}
    [\ce{Ag+}] = 2S, \quad [\ce{X-}] = S, \quad [\ce{Y-}] = S
\end{equation*}
Applying the solubility product:
\begin{equation*}
    K_{sp} = [\ce{Ag+}][\ce{X-}] = (2S)(S) = 2S^2 = 1.0 \times 10^{-12}
\end{equation*}
\begin{equation*}
    S^2 = 0.5 \times 10^{-12} \implies S = \sqrt{0.5 \times 10^{-12}} = \frac{1}{\sqrt{2}} \times 10^{-6} \approx 0.707 \times 10^{-6}\text{ M}
\end{equation*}
The total silver ion concentration is:
\begin{equation*}
    [\ce{Ag+}] = 2S = 2 \times (0.707 \times 10^{-6}) = 1.414 \times 10^{-6}\text{ M} = \sqrt{2} \times 10^{-6}\text{ M}
\end{equation*}

\textbf{Part 2: Limiting Molar Conductivity of $\ce{Y-}$:}
The total conductivity $\kappa_{\text{total}}$ is the sum of the conductivities of all dissolved ions:
\begin{equation*}
    \kappa_{\text{total}} = \kappa(\ce{Ag+}) + \kappa(\ce{X-}) + \kappa(\ce{Y-})
\end{equation*}
Recall that $\kappa_i = \frac{\lambda_i^\circ \cdot c_i}{1000}$:
\begin{equation*}
    \kappa_{\text{total}} = \frac{1}{1000} \left[ \lambda^\circ(\ce{Ag+}) [\ce{Ag+}] + \lambda^\circ(\ce{X-}) [\ce{X-}] + \lambda^\circ(\ce{Y-}) [\ce{Y-}] \right]
\end{equation*}
With $K_{sp}(\ce{AgY}) = 1.0 \times 10^{-10}$ and $K_{sp}(\ce{AgX}) = 1.0 \times 10^{-12}$:
Since $K_{sp}(\ce{AgY}) \gg K_{sp}(\ce{AgX})$, the dominant contributor to $[\ce{Ag+}]$ is $\ce{AgY}$:
\begin{align*}
    [\ce{Ag+}] &\approx [\ce{Y-}] = \sqrt{K_{sp}(\ce{AgY})} = \sqrt{1.0 \times 10^{-10}} = 1.0 \times 10^{-5}\text{ M} \\
    [\ce{X-}] &= \frac{K_{sp}(\ce{AgX})}{[\ce{Ag+}]} = \frac{1.0 \times 10^{-12}}{1.0 \times 10^{-5}} = 1.0 \times 10^{-7}\text{ M}
\end{align*}
Substituting into the total conductivity equation with measured $\kappa_{\text{total}} = 4.04 \times 10^{-6}\ \Omega^{-1}\cdot\text{cm}^{-1}$:
\begin{align*}
    4.04 \times 10^{-6} \times 1000 &= [60.0 \times 1.0 \times 10^{-5}] + [90.0 \times 1.0 \times 10^{-7}] + [\lambda^\circ(\ce{Y-}) \times 1.0 \times 10^{-5}] \\
    4.04 \times 10^{-3} &= 0.60 \times 10^{-3} + 0.009 \times 10^{-3} + \lambda^\circ(\ce{Y-}) \times 1.0 \times 10^{-5} \\
    \lambda^\circ(\ce{Y-}) \times 1.0 \times 10^{-5} &= 3.43 \times 10^{-3} \implies \lambda^\circ(\ce{Y-}) \approx 290.0\ \Omega^{-1}\cdot\text{cm}^2\cdot\text{mol}^{-1}
\end{align*}
This confirms option (a) $290.0$.
\end{solution}

% ==============================================================================
% CHAPTER 1 ANSWER KEY TABLE (FOR MAIN.TEX)
% ==============================================================================
\ifshowsolutions
\else
\newpage
\subsection*{Answer Key: Chapter 1 Practice Problem Bank}
\addcontentsline{toc}{subsection}{Answer Key: Chapter 1}

\begin{center}
\renewcommand{\arraystretch}{1.3}
\begin{tabular}{|c|c||c|c||c|c|}
\hline
\textbf{Problem} & \textbf{Answer} & \textbf{Problem} & \textbf{Answer} & \textbf{Problem} & \textbf{Answer} \\
\hline
1.1 & (b) & 1.7 & (b) & 1.13 & (d) \\
1.2 & (c) & 1.8 & (c) & 1.14 & (a) \\
1.3 & (b) & 1.9 & (d) & 1.15 & (c) \\
1.4 & (a) & 1.10 & (c) & 1.16 & (c) \\
1.5 & (b) & 1.11 & (a), (b), (c) & 1.17 & (a) \\
1.6 & (a) & 1.12 & (b) & 1.18 & Part 1: (b), Part 2: (a) \\
\hline
\end{tabular}
\end{center}
\fi
